% Chapter 2 — Hardware Setup and Commissioning
% Converted from thesis_notes/draft_ch2_hardware.md
\chapter{Hardware Setup and Commissioning}
\label{chap:hardware}
\lhead{\emph{Hardware Setup and Commissioning}}

\section{Overview}
\label{sec:overview}

The experimental plant is a brushed DC gearmotor driven by an L298N H-bridge under pulse-width-modulated (PWM) voltage control, with shaft velocity measured by a quadrature optical encoder. An Arduino Uno R4 Minima generates the PWM command, sets the H-bridge direction, decodes the encoder in hardware interrupts, and streams time-stamped telemetry to a host computer over USB serial. All identification data in this thesis were collected on this single rig; only the firmware behaviour changed between sessions (manual jog versus automated step sequencer, \sref{sec:firmware}).

The signal path is a cascade of four physical stages, and the model identified in later chapters mirrors this same decomposition (\fref{fig:blockdiagram}). The directional voltage asymmetry that motivates the cascade model (\cref{chap:static}) originates in the L298N output stage, while the encoder closes the measurement path.

\begin{figure}[t]
  \centering
  \resizebox{\textwidth}{!}{% Chapter 2 (§2.1) signal-flow / cascade block diagram.
% Package-independent: needs only tikz + libraries
%   positioning, arrows.meta, fit, backgrounds  (all loaded in preamble.tex).
% Usage in a chapter (it is wider than \textwidth, so scale it):
%   \begin{figure}\centering\resizebox{\textwidth}{!}{% Chapter 2 (§2.1) signal-flow / cascade block diagram.
% Package-independent: needs only tikz + libraries
%   positioning, arrows.meta, fit, backgrounds  (all loaded in preamble.tex).
% Usage in a chapter (it is wider than \textwidth, so scale it):
%   \begin{figure}\centering\resizebox{\textwidth}{!}{% Chapter 2 (§2.1) signal-flow / cascade block diagram.
% Package-independent: needs only tikz + libraries
%   positioning, arrows.meta, fit, backgrounds  (all loaded in preamble.tex).
% Usage in a chapter (it is wider than \textwidth, so scale it):
%   \begin{figure}\centering\resizebox{\textwidth}{!}{\input{figures/fig_blockdiagram}}\caption{...}\label{fig:blockdiagram}\end{figure}
\begin{tikzpicture}[
  font=\small,
  >={Latex[length=2mm]},
  block/.style={draw, rounded corners=2pt, minimum height=1.15cm, minimum width=2.35cm,
                align=center, inner sep=3pt, fill=white},
  io/.style={draw, rounded corners=10pt, minimum height=1.15cm, minimum width=1.55cm,
             align=center, inner sep=3pt, fill=black!4},
  al/.style={font=\scriptsize, align=center},
  node distance=12mm and 16mm
]
  \node[io]                          (host)  {Host\\PC};
  \node[block, right=16mm of host]   (mcu)   {Arduino Uno\\R4 Minima};
  \node[block, right=28mm of mcu]    (drv)   {L298N H-bridge\\{\scriptsize\itshape static driver block}};
  \node[block, right=18mm of drv]    (mot)   {DC gearmotor\\{\scriptsize\itshape motor dynamic block}};
  \node[io, right=16mm of mot]       (shaft) {Output\\shaft};
  \node[block, below=18mm of shaft]  (enc)   {600\,PPR encoder\\{\scriptsize 2400 cnt/rev}};

  % forward signal path
  \draw[->] (host) -- node[al, above]{USB} (mcu);
  \draw[->] (mcu)  -- node[al, above]{PWM (D9)\\dir (D4/D5)} (drv);
  \draw[->] (drv)  -- node[al, above]{$V_{\mathrm{motor}}$} (mot);
  \draw[->] (mot)  -- (shaft);

  % feedback path
  \draw[->] (shaft) -- (enc);
  \draw[->] (enc)  -| node[al, pos=0.25, above]{A/B (D2/D3)} (mcu);

  % telemetry return
  \draw[->] (mcu)  to[bend left=22] node[al, below]{telemetry (10\,ms)} (host);

  % cascade-model highlight
  \begin{scope}[on background layer]
    \node[draw, dashed, rounded corners, fit=(drv)(mot), inner sep=11pt,
          label={[al]above:identified cascade model}] {};
  \end{scope}
\end{tikzpicture}
}\caption{...}\label{fig:blockdiagram}\end{figure}
\begin{tikzpicture}[
  font=\small,
  >={Latex[length=2mm]},
  block/.style={draw, rounded corners=2pt, minimum height=1.15cm, minimum width=2.35cm,
                align=center, inner sep=3pt, fill=white},
  io/.style={draw, rounded corners=10pt, minimum height=1.15cm, minimum width=1.55cm,
             align=center, inner sep=3pt, fill=black!4},
  al/.style={font=\scriptsize, align=center},
  node distance=12mm and 16mm
]
  \node[io]                          (host)  {Host\\PC};
  \node[block, right=16mm of host]   (mcu)   {Arduino Uno\\R4 Minima};
  \node[block, right=28mm of mcu]    (drv)   {L298N H-bridge\\{\scriptsize\itshape static driver block}};
  \node[block, right=18mm of drv]    (mot)   {DC gearmotor\\{\scriptsize\itshape motor dynamic block}};
  \node[io, right=16mm of mot]       (shaft) {Output\\shaft};
  \node[block, below=18mm of shaft]  (enc)   {600\,PPR encoder\\{\scriptsize 2400 cnt/rev}};

  % forward signal path
  \draw[->] (host) -- node[al, above]{USB} (mcu);
  \draw[->] (mcu)  -- node[al, above]{PWM (D9)\\dir (D4/D5)} (drv);
  \draw[->] (drv)  -- node[al, above]{$V_{\mathrm{motor}}$} (mot);
  \draw[->] (mot)  -- (shaft);

  % feedback path
  \draw[->] (shaft) -- (enc);
  \draw[->] (enc)  -| node[al, pos=0.25, above]{A/B (D2/D3)} (mcu);

  % telemetry return
  \draw[->] (mcu)  to[bend left=22] node[al, below]{telemetry (10\,ms)} (host);

  % cascade-model highlight
  \begin{scope}[on background layer]
    \node[draw, dashed, rounded corners, fit=(drv)(mot), inner sep=11pt,
          label={[al]above:identified cascade model}] {};
  \end{scope}
\end{tikzpicture}
}\caption{...}\label{fig:blockdiagram}\end{figure}
\begin{tikzpicture}[
  font=\small,
  >={Latex[length=2mm]},
  block/.style={draw, rounded corners=2pt, minimum height=1.15cm, minimum width=2.35cm,
                align=center, inner sep=3pt, fill=white},
  io/.style={draw, rounded corners=10pt, minimum height=1.15cm, minimum width=1.55cm,
             align=center, inner sep=3pt, fill=black!4},
  al/.style={font=\scriptsize, align=center},
  node distance=12mm and 16mm
]
  \node[io]                          (host)  {Host\\PC};
  \node[block, right=16mm of host]   (mcu)   {Arduino Uno\\R4 Minima};
  \node[block, right=28mm of mcu]    (drv)   {L298N H-bridge\\{\scriptsize\itshape static driver block}};
  \node[block, right=18mm of drv]    (mot)   {DC gearmotor\\{\scriptsize\itshape motor dynamic block}};
  \node[io, right=16mm of mot]       (shaft) {Output\\shaft};
  \node[block, below=18mm of shaft]  (enc)   {600\,PPR encoder\\{\scriptsize 2400 cnt/rev}};

  % forward signal path
  \draw[->] (host) -- node[al, above]{USB} (mcu);
  \draw[->] (mcu)  -- node[al, above]{PWM (D9)\\dir (D4/D5)} (drv);
  \draw[->] (drv)  -- node[al, above]{$V_{\mathrm{motor}}$} (mot);
  \draw[->] (mot)  -- (shaft);

  % feedback path
  \draw[->] (shaft) -- (enc);
  \draw[->] (enc)  -| node[al, pos=0.25, above]{A/B (D2/D3)} (mcu);

  % telemetry return
  \draw[->] (mcu)  to[bend left=22] node[al, below]{telemetry (10\,ms)} (host);

  % cascade-model highlight
  \begin{scope}[on background layer]
    \node[draw, dashed, rounded corners, fit=(drv)(mot), inner sep=11pt,
          label={[al]above:identified cascade model}] {};
  \end{scope}
\end{tikzpicture}
}
  \caption{Signal flow and cascade decomposition of the rig. The PWM command passes through the static driver-voltage block (L298N) and the motor dynamic block to produce shaft speed; the 600\,PPR encoder returns the measured speed to the controller, which streams it to the host. The dashed box marks the two blocks identified in \cref{chap:static} and \cref{chap:dynamic}.}
  \label{fig:blockdiagram}
\end{figure}

\section{Components}
\label{sec:components}

The rig comprises the components listed in \tref{tab:components}. Static-characterisation instrumentation (\cref{chap:static}) additionally used a digital multimeter (DMM) for mean motor voltage and a two-channel oscilloscope with a math channel for the differential motor voltage; these form part of the measurement chain.

\begin{table}[t]
  \centering
  \caption{Components of the experimental rig.}
  \label{tab:components}
  \footnotesize
  \begin{tabular}{@{}p{2.4cm}p{3.6cm}p{6.4cm}@{}}
    \toprule
    Subsystem & Part & Key specification \\
    \midrule
    Controller & Arduino Uno R4 Minima & Renesas RA4M1 (Arm Cortex-M4, 48\,MHz) \\
    Motor driver & L298N dual H-bridge (one channel used) & Up to 2\,A/channel; bipolar-transistor output stage with non-negligible saturation drop \\
    Motor & Orange--Johnson 12\,V geared DC motor & 60:1 gearbox, ${\approx}300$\,rpm rated output at 12\,V \\
    Encoder & Orange 3806-OPTI-600-AB-OC & 600\,PPR, 2-channel (A/B) quadrature, open-collector output, on the \textbf{output} shaft \\
    Motor supply & Bench PSU & 12.0\,V set-point, 1.5\,A current limit \\
    Logic supply & Arduino 5\,V rail & feeds encoder $V_{\mathrm{CC}}$ and L298N logic 5\,V \\
    \bottomrule
  \end{tabular}
\end{table}

\section{Wiring and pin mapping}
\label{sec:wiring}

The Arduino-to-peripheral wiring is fixed across all sessions (\tref{tab:pinmap}). Only one L298N channel is used: OUT1/OUT2 drive the motor, ENA/IN1/IN2 are the control inputs, and ENB/IN3/IN4/OUT3/OUT4 are left unconnected. Two jumper changes on the L298N board are required and were verified before every session:
\begin{itemize}
  \item \textbf{ENA jumper removed}, so the Arduino D9 PWM signal, rather than a fixed on-board pull-up, controls channel-A speed.
  \item \textbf{5\,V regulator (5V-EN) jumper removed}, because the Arduino 5\,V rail supplies the L298N logic; leaving it in would back-feed the on-board regulator.
\end{itemize}

The encoder is open-collector, so each channel uses a 4.7\,k$\Omega$ pull-up to the 5\,V rail (placed between the signal line and 5\,V, not in series). All grounds (Arduino, encoder, L298N, and the PSU negative) meet at one star node, and the high-current motor path is kept short and direct rather than routed through thin breadboard rails, to avoid ground-bounce on the shared logic reference.

\begin{table}[t]
  \centering
  \caption{Arduino pin map.}
  \label{tab:pinmap}
  \small
  \begin{tabular}{@{}llp{6cm}@{}}
    \toprule
    Arduino pin & Net & Function \\
    \midrule
    D2  & Encoder A & Quadrature input (CHANGE interrupt) \\
    D3  & Encoder B & Quadrature input (CHANGE interrupt) \\
    D4  & L298N IN1 & Direction bit 1 \\
    D5  & L298N IN2 & Direction bit 2 \\
    D9  & L298N ENA & PWM speed command \\
    5V  & Encoder $V_{\mathrm{CC}}$, L298N logic 5\,V & Logic supply \\
    GND & Encoder GND, L298N GND, PSU $-$ & Common (star) ground \\
    \bottomrule
  \end{tabular}
\end{table}

\section{PWM generation}
\label{sec:pwm}

Speed is commanded as an 8-bit duty value (0--255) on D9 at a fixed \textbf{20\,kHz} carrier. The Uno R4 core does not expose a global PWM-frequency setter, so the carrier is configured directly through the Renesas timer-based PWM peripheral at 20\,kHz, with the duty applied as a percentage of the period. Twenty kilohertz is chosen deliberately: it sits at the top of the audible band so the motor does not whine, and, more importantly for this study, it makes the H-bridge output a clean high-frequency square wave whose \textbf{time-average} is a stable, well-defined motor voltage. That averaging is what allows the DMM and oscilloscope math channel to read a repeatable $V_{\mathrm{motor}}$ during the static sweeps of \cref{chap:static}; at a low carrier the average reading wanders and the static map becomes unmeasurable.

\section{Encoder interface and velocity estimation}
\label{sec:encoder}

The encoder produces 600 pulses per revolution per channel. Both channels are decoded in full quadrature: A and B each trigger an interrupt on every logic edge, and a 16-entry state-transition lookup table converts each (previous, current) two-bit state pair into a step of $-1$, $0$, or $+1$. Four edges per pulse period yield \textbf{2400 counts per output-shaft revolution}, i.e.\ an angular resolution of $360^\circ/2400 = \mathbf{0.15^\circ}$ per count at the output shaft.

Velocity is not measured directly; it is estimated by finite difference of the accumulated count over the fixed 10\,ms telemetry interval,
\begin{equation}
  \mathrm{RPM} = \frac{\Delta\mathrm{count}\times 60\times 10^{6}}{2400\times \Delta t_{\mu\mathrm{s}}},
  \label{eq:rpm}
\end{equation}
evaluated once per sample. A single-count change over a 10\,ms window therefore corresponds to a velocity quantisation of $60\times 10^{6}/(2400\times 10^{4}) = 2.5$\,rpm/count, which sets the noise floor of the RPM signal and is the reason the dynamic-fitting code (\cref{chap:dynamic}) treats sub-${\approx}2$\,rpm excursions as quantisation rather than motion. The encoder ISR is short and non-blocking; count reads in the main loop are taken inside a brief interrupt-disabled critical section to avoid tearing the 32-bit counter.

\section{Direction and sign conventions}
\label{sec:conventions}

Direction is set by the IN1/IN2 pair: forward drives IN1 high / IN2 low, reverse drives IN1 low / IN2 high, and IN1 low / IN2 low coasts. The convention adopted throughout the thesis is \textbf{forward $\to$ positive RPM, reverse $\to$ negative RPM}. Each telemetry row carries the direction explicitly as a single character (F for forward, R for reverse, N for neutral/coast), so that the unsigned PWM magnitude and the signed velocity can always be reconstructed as $\mathrm{signed} = \mathrm{pwm}\times s$, with $s = +1, -1, 0$ for forward, reverse, and neutral respectively.

\section{Firmware architecture}
\label{sec:firmware}

Two firmware behaviours were used, sharing identical pin definitions, encoder decoding, PWM setup, and telemetry format, and differing only in how transitions are commanded:
\begin{itemize}
  \item \textbf{Manual-mode firmware} accepts single-character serial commands (set forward / reverse / stop, zero the counter, print help) and an integer PWM magnitude, and streams telemetry continuously. It was used for the operator-driven static sweeps and deadzone ramps of \cref{chap:static}.
  \item \textbf{Automated step-sequencer firmware} (a distinct trial array for the Phase~2.1 and gap-fill sessions) executes a pre-programmed array of step trials without operator intervention on a single start command, emitting each trial as a block bracketed by explicit start and end markers. This removed operator timing variability from the dynamic step-response data of \cref{chap:dynamic}.
\end{itemize}

All telemetry is emitted at \textbf{230400 baud} as comma-separated rows of time stamp, PWM command, direction, encoder count, and velocity at a 10\,ms cadence. A host-side capture script logs the stream, and a splitter script cuts the delimited blocks into per-trial CSV files; raw CSVs are treated as immutable and all derived quantities are written separately.

\section{Commissioning and acceptance (Phase 0)}
\label{sec:commissioning}

Before any identification data were trusted, the rig passed an acceptance gate confirming a clean firmware build, the 12.0\,V / 1.5\,A supply setting, and the 20\,kHz PWM command. The substantive commissioning finding concerned the \textbf{encoder sign}. A single-operating-point check was run at PWM~=~150 in both directions, with oscilloscope CH1 on OUT1, CH2 on OUT2, and a math channel computing CH1~$-$~CH2 as the differential motor voltage (\tref{tab:scopecheck}).

\begin{table}[t]
  \centering
  \caption{Phase~0 single-point scope check at PWM~=~150.}
  \label{tab:scopecheck}
  \begin{tabular}{@{}lccc@{}}
    \toprule
    Command & Math (CH1$-$CH2) mean & Serial RPM & PSU current \\
    \midrule
    Forward & $+1.7$ to $+2.1$\,V & $-58$ to $-59$\,rpm & 0.22\,A \\
    Reverse & $-2.5$ to $-3.0$\,V & $+56$ to $+61$\,rpm & 0.21\,A \\
    \bottomrule
  \end{tabular}
\end{table}

The motor-voltage polarity was internally consistent (a forward command produced a positive differential voltage and a reverse command a negative one), but the \textbf{reported RPM sign was inverted relative to the command}: forward drive read negative, reverse read positive. The diagnosis was that the encoder A/B phasing was reversed with respect to the chosen command convention, not a wiring or polarity fault in the power path. The fix was a one-line firmware change setting the encoder decode sign to $-1$, after which the build was rebuilt successfully and the PWM~=~150 gate was repeated to confirm forward $\to$ positive RPM and reverse $\to$ negative RPM. This sign correction is baked into every firmware variant used for the data in this thesis, and it establishes the directional convention used uniformly from \cref{chap:static} onward.

A separate within-session repeatability spot-check at PWM~=~$\pm 255$ (morning versus evening of the same day) confirmed the measurement chain was stable: the motor-side velocity reproduced to within ${\approx}1\,\%$, while the driver-side motor voltage and PSU current drifted by a few percent, an early and direct indication that what little instability the rig exhibits lives in the L298N (driver) stage rather than the motor, foreshadowing the cascade decomposition developed in \cref{chap:static}.

\section{Summary}
\label{sec:hw-summary}

The plant is a 12\,V, 60:1 DC gearmotor driven by one channel of an L298N at a 20\,kHz, 8-bit PWM command, with direction set by two logic lines and velocity measured by a 600\,PPR quadrature encoder decoded to 2400 counts/rev ($0.15^\circ$/count) on the output shaft. An Arduino Uno R4 Minima performs PWM generation, direction control, interrupt-driven encoder decoding, and 10\,ms telemetry at 230400 baud; velocity is a finite-difference estimate with a 2.5\,rpm quantisation floor. Commissioning fixed an inverted encoder sign and established the forward-positive / reverse-negative convention, and an early repeatability check localised the rig's only measurable drift to the driver stage. This hardware is the fixed substrate for the static characterisation (\cref{chap:static}), dynamic characterisation (\cref{chap:dynamic}), and validation (\cref{chap:validation}) that follow.
